% New mathematical fragment. Uses only standard amsmath/amsthm notation.
% The enclosing report supplies theorem, lemma, corollary environments.
\section{Equality from one fixed probe of a bound witness}
\label{sec:inferred-discharge}

This section concerns a new restriction of the equality-bearing discharge
rule. Its contribution is an exact reduction of the rule's proof interface:
the equality premise and its supplied right-hand term are both removed.
Their replacement is a fixed syntactic probe of the already proved major
formula. The theorem is an identity-language equivalence, rather than merely
an equiconsistency result. No assertion of absolute minimality follows from
it. Whether its complete implementation is shorter is a separate measured
question.

\subsection{Syntax, ordinary reading, and all primitive rules}

There is one numerical sort. Begin with raw terms
\[
 t::=x\mid0\mid S(t)\mid I(n,a;u.h).
\]
Only $u$ in $h$ is bound. Write $N$ for the full normal-form operation for
the context-closed rules
\[
 I(0,a;u.h)\longrightarrow a,
 \qquad I(Sn,a;u.h)\longrightarrow h[I(n,a;u.h)/u].
\]
The terms used in the calculus are normal forms. All displayed substituted
instances are capture avoiding and followed by $N$. Normalization is
metasyntactic and is included in the proof checker; it is not an additional
object-language axiom. On formulas, $N$ acts on terms and preserves all
logical constructors. The normalization facts used here are termination,
confluence, and the substitution coherence identity
\[
 N(N(a)[\sigma])=N(a[\sigma]),
\]
with simultaneous capture-avoiding substitution on either side. These facts
belong to the inherited full-normalization comparison; they are not replaced
by a numerical claim that a universal primitive-recursive evaluator exists.

Formulas have the grammar
\[
 A::=(s=t)\mid D_{x,y}B(x,y),
\]
where $D$ binds both displayed numerical variables. A clause has the form
$[\Delta;y.B(y)]$, where $\Delta$ is either empty or consists of one
formula $A$, and $y$ does not occur free in $\Delta$. Its ordinary reading is
\[
 [\,;y.B]^\circ=\forall y\,B^\circ(y),\qquad
 [[A];y.B]^\circ=A^\circ\to\forall y\,B^\circ(y).
\]
The ordinary readings of atoms and $D$ are
\[
 (s=t)^\circ=(s=t),\qquad
 (D_{x,y}B)^\circ=\exists x\bigl(B^\circ(x,0)\to
                                  \forall y\,B^\circ(x,y)\bigr).
\]
A context is a finite list of clause occurrences. An ordinary assumption
$A$ abbreviates the unconditional constant clause $[\,;v.A]$ with unused
$v$. Every context contains the distinguished fixed axiom occurrence
\[
 \mathcal R=[\,;v.v=v].
 \tag{Ax}
\]
In particular, reflexivity is an explicit axiom clause, rather than a theorem
obtained with no equality assumption of any kind. Put $H=(0=0)$ and
$1=S0$, $2=S(S0)$.

The first primitive rule is clause use:
\[
 \frac{\{\,\Gamma\vdash A\mid A\text{ is the antecedent occurrence in }
                                  \Delta\,\}}
      {\Gamma\vdash N(B(t))}
 \quad [\Delta;y.B]\in\Gamma.
 \tag{U}
\]
Thus an empty antecedent requires no premise and a singleton antecedent
requires exactly one proof. In particular a use of $\mathcal R$ gives $t=t$.

The second primitive rule is paired introduction:
\[
 \frac{\Gamma,B(t,\eta)\vdash B(t,Sz)}
      {\Gamma\vdash D_{x,y}B(x,y)},\qquad
 \eta=\begin{cases}
 z,&B\text{ has equation root and }t\text{ is literally }0,\\
 0,&\text{otherwise}.
 \end{cases}
 \tag{I}
\]
Here $z$ is fresh for $\Gamma,t$, and all passive parameters of $B$.
The witness $t$ is canonical, so the zero test is a test of its canonical
syntax. The local assumption is an unconditional constant clause.

For the final primitive rule, suppose a proved major is $D_{x,y}B(x,y)$
and choose a fresh numerical parameter $b$. Let
\[
 B_b(y)=N(B(b,y)),\qquad
 \mathcal B_b=[[B_b(0)];y.B_b(y)].
\]
Define a term determined by this major and the fresh name:
\[
 \tau_B(b)=
 \begin{cases}
 r,&B_b(0)\equiv(0=0)\text{ and }B_b(1)\equiv(b=r),\\
 b,&\text{otherwise}.
 \end{cases}
 \tag{P}
\]
The symbol $\equiv$ means literal equality of canonical syntax. In the
first branch the right-hand term $r$ is read from the equation; it is not
supplied independently. The asymmetry of this probe is intentional. No test
of provable equality, no search for an equality proof, and no arbitrary
choice of specialization target is involved.

The inferred-discharge rule is
\[
 \frac{\Gamma\vdash D_{x,y}B(x,y)\qquad
       \Gamma,\mathcal B_b\vdash F}
      {\Gamma\vdash N(F[\tau_B(b)/b])}.
 \tag{E$^*$}
\]
The parameter $b$ is fresh for $\Gamma$ and all passive parameters of the
major. The exported normal form must be free of $b$. There is no separate
equality premise. The synthesized term $\tau_B(b)$ is allowed to contain
$b$; the substitution is simultaneous, so an inserted copy is not traversed
again. These three rules and Ax are the complete primitive inventory.

\begin{lemma}[Soundness of inferred discharge]
\label{lem:inferred-sound}
Rule E$^*$ is sound in intuitionistic first-order equality logic over the
scalar computation equations, with the ordinary readings just specified.
\end{lemma}
\begin{proof}
Use existential elimination on the major at $b$. Its local conditional
universal assertion is exactly $\mathcal B_b^\circ$. If the first branch
of P applies, reflexivity proves its base $H$ and specialization at $1$
gives $b=r$. If P chooses its fallback, reflexivity gives $b=b$ instead.
Thus, in both cases, the opened context proves $b=\tau_B(b)$.
The second premise gives $F$, so equality substitution gives
$F[\tau_B(b)/b]$ and computation identifies this with its normal form.
The stated absence of $b$ from that normal form licenses existential
elimination. Equality substitution remains valid for $b=r(b)$; it is not
recursive replacement of the inserted term.
\end{proof}

Rule U has its ordinary clause reading. Rule I is sound by numerical
coverage in the case branch and by induction for literal equations in the
progress branch. Accordingly all primitive rules are sound in the same
ordinary scalar theory as the inherited clause calculus.

\subsection{Structural operations and cases before equality}

We first recover the tools needed to derive equality. None of the following
arguments presupposes equality transport or arbitrary numerical substitution
in derivations.

\begin{lemma}[Ordinary paired elimination]
\label{lem:inferred-ordinary-elim}
The rule
\[
 \frac{\Gamma\vdash D_{x,y}B(x,y)\qquad
       \Gamma,\mathcal B_b\vdash A}
      {\Gamma\vdash A},\qquad b\notin\operatorname{FV}(A),
\]
is derivable using one E$^*$.
\end{lemma}
\begin{proof}
Use the displayed branch proof with $F=A$. Regardless of which branch of
P is selected, $N(A[\tau_B(b)/b])=A$ literally: $A$ is canonical and contains
no free $b$. This argument also works when the synthesized term contains
$b$, and does not require any control of the probe.
\end{proof}

\begin{lemma}[Structural operations needed below]
\label{lem:inferred-structural}
Weakening by clause occurrences, injective renaming of free numerical
parameters, and cut of an ordinary formula assumption are admissible.
\end{lemma}
\begin{proof}
For weakening, preserve the selected occurrence at every U step and freshen
the finitely many eigenparameters in I and E$^*$ as necessary. Induct on
the proof tree. For injective renaming, use equivariance of normalization
and capture-avoiding substitution. The fresh parameter remains a fresh
parameter, $0$ and $1$ are unchanged, and injectivity preserves and reflects
the literal pattern defining P. The literal-zero witness test in I is also
preserved by injective renaming.

For ordinary-assumption cut, replace a U step selecting the designated
unconditional constant clause by the supplied proof, weakened into its
current local context. All other selected clause occurrences are retained,
even if they have the same value as the cut occurrence. Recurse through I
and E$^*$ after freshening their eigenparameters. These constructions only
insert derivations and rename parameters; they do not substitute an
arbitrary numerical term into an existing derivation.
\end{proof}

For an unused $w$, write $q_yA(y)=D_{w,y}A(y)$.
\begin{lemma}[Closure introduction, specialization, and numerical cases]
\label{lem:inferred-cases}
The following rules are derivable before equality transport:
\[
 \frac{\Gamma,A(0)\vdash A(Sz)}{\Gamma\vdash q_yA(y)},
 \qquad
 \frac{\Gamma\vdash q_yA(y)\qquad\Gamma\vdash A(0)}
      {\Gamma\vdash A(t)}.
\]
Consequently, if $\Gamma\vdash A(0)$ and $\Gamma\vdash A(Sz)$ with $z$
fresh, then $\Gamma\vdash A(t)$ for every well-scoped term $t$.
\end{lemma}
\begin{proof}
Introduce the unused witness coordinate of $q$ at $1$. Since this witness
is not literally zero, I uses its case branch at every matrix root.
For specialization, use Lemma~\ref{lem:inferred-ordinary-elim}; its opened
clause is $[[A(0)];y.A(y)]$ because $w$ is unused. Apply U with the given
base proof and target $t$. The result is independent of the fresh witness,
so ordinary elimination closes. For numerical cases, weaken the supplied
successor proof by the unused assumption $A(0)$ and compose the two rules.
\end{proof}

\subsection{All equality from one probe}

The selector and parity terms used here are abbreviations in the numerical
language:
\[
 C(v,a,c)=N(I(v,a;u.c)),\quad u\notin\operatorname{FV}(c),
 \qquad P(y)=N(I(y,0;u.C(u,1,0))).
\]
Computation gives the literal identities
\[
 C(0,a,c)\equiv a,\quad C(Sv,a,c)\equiv c,
 \qquad P(1)\equiv1,\quad P(2)\equiv0.
\]
No theorem about the range of $P$ and no induction for $P$ is required.

\begin{theorem}[Uncrossed transport from inferred discharge]
\label{thm:inferred-transport}
For every numerical template $h(w)$, the rule
\[
 \frac{\Gamma\vdash s=t\qquad\Gamma\vdash N(h(s))=r}
      {\Gamma\vdash N(h(t))=r}
\]
is derivable. Apart from the ordinary eliminations used to implement
numerical cases, it uses one E$^*$ with a nontrivial synthesized name.
\end{theorem}
\begin{proof}
Choose fresh $x,y$ and define
\[
 B(x,y)=
 \bigl(C(y,0,C(P(y),N(h(x)),x))
      =C(y,0,C(P(y),r,t))\bigr).
\]
Its three relevant rows are literal normalized formulas:
\[
 B(x,0)\equiv H,\qquad B(x,1)\equiv(x=t),\qquad
 B(x,2)\equiv(N(h(x))=r).
 \tag{R}
\]
At witness $s$, its successor row is
\[
 C(P(Sz),N(h(s)),s)=C(P(Sz),r,t).
\]
Apply numerical cases from Lemma~\ref{lem:inferred-cases} to
\[
 A(v)=\bigl(C(v,N(h(s)),s)=C(v,r,t)\bigr)
\]
at the target $P(Sz)$. The zero case is the supplied $N(h(s))=r$;
every successor case is the supplied $s=t$. Choose the numerical-case
eigenparameter independently of $z$ and fresh for its entire local context.
This derives $B(s,Sz)$ without using I's temporary assumption. Hence I
introduces $D_{x,y}B(x,y)$ at $s$ even if $s$ is literally zero and the
progress branch, rather than the case branch, is selected.

Apply E$^*$ to that major at a fresh $b$. By R, the probe synthesizes
exactly $\tau_B(b)=t$, and the opened clause has base $H$. Use Ax at zero
and then U at $2$ to derive the branch formula $N(h(b))=r$.
The exported formula is
\[
 N\bigl((N(h(b))=r)[t/b]\bigr)=(N(h(t))=r),
\]
by normalization--substitution coherence. It is free of $b$ because the
passive data $h,t,r$ were chosen before $b$. All uses of numerical cases
above were available without equality, so the construction is not circular.
\end{proof}

\begin{corollary}[Ordinary equational reasoning]
\label{cor:inferred-equality}
Symmetry, transitivity, term congruence, replacement in either side of an
equation, and crossed transport are derivable.
\end{corollary}
\begin{proof}
For symmetry, in Theorem~\ref{thm:inferred-transport} use $h(w)=w$ and
the reflexive minor $s=s$; the result is $t=s$. Applying the theorem to
$s=t$ and $h(s)=r$, then using symmetry, gives crossed transport
$r=h(t)$. For transitivity from $s=t$ and $t=u$, first reverse $s=t$ and
apply uncrossed transport with identity template to obtain $s=u$.
Term congruence follows from a reflexive minor $h(s)=h(s)$ and symmetry
as necessary. For a two-sided equation $a(s)=b(s)$, uncrossed transport
first gives $a(t)=b(s)$; symmetry, a second use with template $b$, and
symmetry give $a(t)=b(t)$. All terms in these statements are understood in
normal form.
\end{proof}

\begin{theorem}[Transport in every formula template]
\label{thm:inferred-formula-transport}
For every fixed formula template $F(v)$,
\[
 \Gamma\vdash s=t,\quad\Gamma\vdash N(F(s))
 \quad\Longrightarrow\quad\Gamma\vdash N(F(t)).
\]
\end{theorem}
\begin{proof}
Induct externally on formula depth. The equation case is
Corollary~\ref{cor:inferred-equality}. In the compound case write
$F(v)=D_{x,y}B(v,x,y)$. Open the supplied $D_{x,y}B(s,x,y)$ by ordinary
paired elimination at a fresh $c$. Introduce $D_{x,y}B(t,x,y)$ at witness
$c$. A fresh parameter is not literally zero, so I uses its case branch.
Its local assumption is $B(t,c,0)$. By symmetry and the induction
hypothesis, obtain $B(s,c,0)$. Specialize the old opened clause at $Sz$,
giving $B(s,c,Sz)$, and apply the induction hypothesis forward to obtain
$B(t,c,Sz)$. Close I and ordinary elimination. Every recursive equality
transport concerns the strictly smaller matrix $B$. The argument is an
induction on finite syntax in the metatheory, not an added induction schema
for quantified formulas in the object theory.
\end{proof}

\subsection{Exact comparison and the iPRA interpretation boundary}

Let $\mathcal E$ be the inherited canonical clause calculus with Ax, U,
I, and equality-bearing discharge
\[
 \frac{\Gamma\vdash D_{x,y}B(x,y)\quad
       \Gamma,\mathcal B_b\vdash b=t\quad
       \Gamma,\mathcal B_b\vdash F}
      {\Gamma\vdash N(F[t/b])},
 \tag{E$^=$}
\]
with exactly the same freshness and exported-scope conditions as E$^*$.
Write $\mathcal N$ for the three-rule calculus defined here.

\begin{theorem}[Identity-language equivalence]
\label{thm:inferred-equivalence}
For every finite canonical clause context $\Gamma$ containing Ax and every
canonical formula $A$,
\[
 \Gamma\vdash_{\mathcal N}A
 \quad\Longleftrightarrow\quad
 \Gamma\vdash_{\mathcal E}A.
\]
The translations are effective transformations of finite proof trees and
preserve the entire end sequent literally. The theorem applies to arbitrary
retained clauses, not only to those reachable in a closed proof.
\end{theorem}
\begin{proof}
Forwards, translate recursively. In an E$^*$ instance, if P selects its
first branch, use Ax and U at $1$ in the opened context to prove
$b=\tau_B(b)$. Otherwise use Ax to prove $b=b$. The translated major,
this equality proof, and the translated branch are premises of E$^=$.
The resulting conclusion and scope condition are unchanged. U and I are
identical in the two systems.

Backwards, translate recursively. In an E$^=$ instance, open the translated
major using the derived ordinary elimination of
Lemma~\ref{lem:inferred-ordinary-elim}. The translated branch proofs give
$b=t$ and $F$ in that opened context. Apply
Theorem~\ref{thm:inferred-formula-transport}, with the numerical variable $b$
as source, to obtain $N(F[t/b])$. The stipulated absence of $b$ permits
ordinary elimination. The equality theorem was proved within $\mathcal N$
without using the present translation, so this recursion on the original
proof tree is well founded. Freshen eigenparameters before inserting the
finite replacement derivations.
\end{proof}

\begin{corollary}[Bi-interpretability and equiconsistency]
\label{cor:inferred-ipra}
Composed with the inherited comparison of $\mathcal E$ with the ordinary
scalar theory and of that scalar theory with full first-order iPRA,
$\mathcal N$ is equiconsistent and bi-interpretable with full first-order
intuitionistic primitive recursive arithmetic. The numerical domains are
whole, equality is actual equality, and both numerical composite
isomorphisms have graph $x=y$.
\end{corollary}
\begin{proof}
The new equivalence is the identity on all terms, formulas, and clause
meanings. It therefore composes with the inherited interpretations without
altering either numerical map or any formula translation. The new legs of
each round trip are literally identity translations; the inherited legs
are provably equivalent to their inputs and have the identity numerical
graph. This gives the two composites, rather than merely interpretations
in opposite directions. The designated contradiction $0=S0$ is unchanged,
so equiconsistency follows as a special case.
\end{proof}

The last corollary has a precise dependency boundary. The new work is
Theorem~\ref{thm:inferred-equivalence}; it does not reprove the inherited
scalar construction of predecessor, equality reflection, pairing, and a
compiler for all primitive-recursive definitions. The report must cite the
source snapshot and those comparison results. Existing Lean results for the
scalar and native systems do not by themselves constitute a Lean
verification of E$^*$ or of its Python checker.

\paragraph{Arbitrary substitution is an admissible operation, not an
unchecked rule transformation.}
An arbitrary numerical substitution can create the literal probe pattern
in P, destroy it, or turn I's witness into zero. Thus one must not claim
that substitution preserves primitive rule tags by a direct induction.
The structural operations actually used in the proof above are only
weakening, injective fresh renaming, and ordinary-assumption cut, for which
the direct arguments were supplied. General substitution, when wanted,
is recovered by translation to an inherited substitution-stable calculus,
applying its substitution theorem, and translating back. This procedure
also establishes admissibility when a synthesized name contains its own
fresh parameter. No stronger claim about primitive-recursive bounds for
full normalization or for all these proof translations is made.
